\onecolumn
%% \begin{longtable}{|p|p|p|p|p|p|p|}
%% \begin{longtable}{|p{1.5cm}|p{1cm}|p{1.5cm}|p{1cm}|p{3.5cm}|p{2.5cm}|p{11cm}|}
\begin{landscape}
    \setlength\LTleft{0pt}
    \setlength\LTright{0pt}
    \begin{longtable}{|p{1.5cm}|l|p{1.5cm}|p{1.3cm}|p{2.75cm}|p{2.5cm}|p{7.5cm}|p{4.75cm}|}
        \tabcap{Study Descriptor table}\label{tab:descriptor-table} \\
        \hline
        \textbf{Author(s)} & \textbf{Design} & \textbf{Population} & \textbf{Mean Age~(yrs.)} & \textbf{Intervention} & \textbf{Comparator} & \textbf{Main Results} & \textbf{Safety} \\
        \hline
        \endfirsthead

        \multicolumn{8}{c}%
        {{\bfseries \tablename\ \thetable{} -- continued from previous page}} \\
        \hline
        \textbf{Author(s)} & \textbf{Design} & \textbf{Population} & \textbf{Mean Age (yrs.)} & \textbf{Intervention} & \textbf{Comparator} & \textbf{Main Results} & \textbf{Safety} \\
        \hline
        \endhead

        \hline
        \multicolumn{8}{|r|}{Continued on next page} \\ 
        \hline
        \endfoot

        \hline
        \multicolumn{8}{|p{23cm}|}{\small Abbreviations: AE=Adverse Events, BCVA=Best Corrected Visual Acuity, CI=Confidence Interval, CMT=Central Macular Thickness, CNV=Choroidal Neovascularization, CRT=Central Retinal Thickness, ITT=Intent-to-Treat, IVR=Intravitreal Ranibizumab, NR=Not Reported, PRN=\PRN, RCT=Randomized Controlled Trial, SAE=Seriouos Adverse Events, SD=Stanard Deviation, T\&E/TREX=Treat-and-Extend, T\&O=Treat \& Observe, W\&E=Wait \& Extend, yrs.=Years} \\
        \hline
        \endlastfoot
        
        \citet{catt2012} & RCT & \num{1107}, \qty{38.12}{\percent}~male & 79.27 & Ranibizumab + \mbox{Bevacizumab} & Placebo & Over two years, both drugs showed similar average visual acuity \mbox{improvements}; Bevacizumab-Ranibizumab varied by 1.4 letters (\qty{95}{\percent} CI,~0.8-3.7;~P=0.21). Monthly~\mbox{treatment} led to a greater mean improvement than PRN, by 2.4~\mbox{letters} (\qty{95}{\percent}~CI,~0.1-4.8;~P=0.046). & No significant difference in death or arteriothrombotic events. Higher rate of systemic SAEs with Bevacizumab.\\ \hline
        \citet{eldem2015} & RCT & \num{93}, \qty{53.76}{\percent}~male & 70.35 & Ranibizumab & Placebo & The efficacy of both T\&O and W\&E Ranibizumab regimens showed comparable improvements in BCVA for AMD patients. & Similar AEs in both treatment groups, mostly mild or transient. No SAEs \mbox{related} to Ranibizumab.\\ \hline
        \citet{ho2014} & RCT & \num{1098} & NR & 0.5,~or~2.0~mg monthly IVR or PRN after 3~monthly loading doses & Placebo & At month 24, the mean change from baseline in BCVA was (letters) +9.1~(0.5~mg~monthly), +7.9~(0.5~mg~PRN), +8.0~(2.0~mg~monthly), and +7.6 (2.0~mg PRN). & Similar safety profile among all groups. Comparable ocular and \mbox{systemic} AEs, no new safety \mbox{concerns}.\\ \hline
        \citet{kertes2019} & RCT & \num{526}, \qty{39.7}{\percent}~male & 78.8 & Ranibizumab (T\&E) & Monthly regimen of Ranibizumab & Mean BCVA improvement of 8.4~letters (SD,~11.9~\mbox{letters}) and 6.0~letters (SD,~11.9~\mbox{letters};~P=0.017) in the T\&E and monthly regimens, respectively. & Well-tolerated treatment. AEs, SAEs, and ophthalmic examinations did not show new safety concerns.\\ \hline
        \citet{kertes2020} & RCT & \num{580}, \qty{39.7}{\percent}~male & 78.8 & Ranibizumab (T\&E) & Monthly regimen of Ranibizumab & The mean (SD) improvement in BCVA was comparable \mbox{between} the T\&E approach and the monthly approach, with values of 6.8 (14.1) letters and 6.0 (12.6) letters, respectively. & Similar incidence of AEs between T\&E and monthly treatment groups. No new safety concerns identified.\\ \hline
        \citet{nunes2019} & RCT & \num{45}, \qty{46.67}{\percent}~male & 75.29 & Bevacizumab monthly, or \mbox{biweekly}, or Ranibizumab monthly & Placebo & Statistically significant improvements in visual acuity and reduction in CMT were observed in all groups. Nevertheless, no significant differences were found among the three groups. & Similar AEs across groups. Mild conjunctival hyperemia, subconjunctival hemorrhage, transient floaters. No \mbox{serious} local complications.\\ \hline
        \citet{russo2018} & RCT & \num{58}, \qty{43.10}{\percent}~male & 75.05 & Ribizumab and \mbox{Indomethacin} & Ranibizumab Monotherapy & After 12 months, notable enhancements were observed in both groups regarding BCVA and CRT. & No SAEs, similar occurrence of ocular AEs in both groups. Most common was mild burning or stinging.\\ \hline
        \citet{schmidt2011} & RCT & \num{293}, 40.79\%~male & 79.3 & 0.3, or \qty{0.5}{\milli\gram}~quarterly, or \qty{0.3}{\milli\gram}~monthly doses of Ranibizumab & Placebo & The mean decrease in CRT from baseline to month 12 in the ITT population was \qty{-96.0}{\micro\metre} in \qty{0.3}{\milli\gram} quarterly, \qty{-105.6}{\micro\metre} in 0.5~mg quarterly, and \qty{-105.3}{\micro\metre} in \qty{0.3}{\milli\gram} monthly group. & Similar incidences of AEs. \mbox{Common} ocular AEs: eye pain, \mbox{conjunctival} hemorrhage. Non-ocular AEs: \mbox{nasopharyngitis}, hyper\-tension.\\ \hline
        \citet{semeraro2015} & RCT & \num{75}, \qty{48}{\percent}~male & 76.7 & Ketorolac eyedrops~+~IVR, or Verteporfin photodynamic~therapy~+~IVR & IVR Monotherapy & Each group exhibited noteworthy enhancement in both BCVA and CRT. The Ranibizumab plus Ketorolac and Ranibizumab with Verteporfin groups necessitated a lower number of IVR treatments compared to the Ranibizumab Monotherapy group. & Comparable safety profile across treatment groups. Mild to moderate ocular AEs, no SAEs noted.\\ \hline
        \citet{wykoff2018} & RCT & \num{60} & NR & Ranibizumab (TREX) & Monthly check-ups, treated if exudative disease activity & The shift from a monthly schedule to a PRN approach was linked to a decrease in BCVA. & Small proportion experienced ocular SAEs. No instances of endoph\-thalmitis or intra\-ocular inflammation.\\ \hline
        \end{longtable}
        \twocolumn
        %% \noindent
\end{landscape}